\documentclass[10pt,a4paper]{amsart}
\usepackage{amsmath,amssymb,amsthm,graphicx}
\usepackage{verbatim}
\newenvironment{reference}{\comment}{\endcomment}
%\newenvironment{reference}{}{}
\newenvironment{slogan}{\comment}{\endcomment}
\newenvironment{history}{\comment}{\endcomment}

% For commutative diagrams we use Xy-pic
\usepackage[all]{xy}

% We use 2cell for 2-commutative diagrams.
\xyoption{2cell}
% Theorem environments.
%
\theoremstyle{plain}
\newtheorem{theorem}[subsection]{Theorem}
\newtheorem{proposition}[subsection]{Proposition}
\newtheorem{lemma}[subsection]{Lemma}

\theoremstyle{definition}
\newtheorem{definition}[subsection]{Definition}
\newtheorem{example}[subsection]{Example}
\newtheorem{exercise}[subsection]{Exercise}
\newtheorem{situation}[subsection]{Situation}

\theoremstyle{remark}
\newtheorem{remark}[subsection]{Remark}
\newtheorem{remarks}[subsection]{Remarks}

\numberwithin{equation}{subsection}

% Macros
%
\def\lim{\mathop{\mathrm{lim}}\nolimits}
\def\colim{\mathop{\mathrm{colim}}\nolimits}
\def\Spec{\mathop{\mathrm{Spec}}}
\def\Hom{\mathop{\mathrm{Hom}}\nolimits}
\def\Ext{\mathop{\mathrm{Ext}}\nolimits}
\def\SheafHom{\mathop{\mathcal{H}\!\mathit{om}}\nolimits}
\def\SheafExt{\mathop{\mathcal{E}\!\mathit{xt}}\nolimits}
\def\Sch{\mathit{Sch}}
\def\Mor{\mathop{\mathrm{Mor}}\nolimits}
\def\Ob{\mathop{\mathrm{Ob}}\nolimits}
\def\Sh{\mathop{\mathit{Sh}}\nolimits}
\def\NL{\mathop{N\!L}\nolimits}
\def\CH{\mathop{\mathrm{CH}}\nolimits}
\def\proetale{{pro\text{-}\acute{e}tale}}
\def\etale{{\acute{e}tale}}
\def\QCoh{\mathit{QCoh}}
\def\Ker{\mathop{\mathrm{Ker}}}
\def\Im{\mathop{\mathrm{Im}}}
\def\Coker{\mathop{\mathrm{Coker}}}
\def\Coim{\mathop{\mathrm{Coim}}}

% Boxtimes
%
\DeclareMathSymbol{\boxtimes}{\mathbin}{AMSa}{"02}

%
% Macros for moduli stacks/spaces
%
\def\QCohstack{\mathcal{QC}\!\mathit{oh}}
\def\Cohstack{\mathcal{C}\!\mathit{oh}}
\def\Spacesstack{\mathcal{S}\!\mathit{paces}}
\def\Quotfunctor{\mathrm{Quot}}
\def\Hilbfunctor{\mathrm{Hilb}}
\def\Curvesstack{\mathcal{C}\!\mathit{urves}}
\def\Polarizedstack{\mathcal{P}\!\mathit{olarized}}
\def\Complexesstack{\mathcal{C}\!\mathit{omplexes}}
% \Pic is the operator that assigns to X its picard group, usage \Pic(X)
% \Picardstack_{X/B} denotes the Picard stack of X over B
% \Picardfunctor_{X/B} denotes the Picard functor of X over B
\def\Pic{\mathop{\mathrm{Pic}}\nolimits}
\def\Picardstack{\mathcal{P}\!\mathit{ic}}
\def\Picardfunctor{\mathrm{Pic}}
\def\Deformationcategory{\mathcal{D}\!\mathit{ef}}
\title{Degrees of field extensions are multiplicative in towers}
\author{The Stacks Project Authors}
\date{stacks-project:fields.tex:local-snapshot-20261012; GNU FDL 1.2 or later}
\begin{document}
\maketitle
\noindent\textit{Editorial note.} All delivered mathematical text is the author's own, copied from the pinned Stacks Project chapter file fields.tex; no statement or proof step was rewritten, reordered or removed.\par\medskip
\noindent\textit{Editorial note.} The counted proof is the author's complete argument: This is the workhorse behind every degree count in the chapter, including the multiplicativity steps invoked by the embedding-count lemmas.\par\medskip
\noindent\textit{Editorial note.} The label lemma-multiplicativity-degrees names the counted statement, and the proof environment immediately following it in the source is delivered verbatim as its proof.\par\medskip
\noindent\textit{Editorial note.} The counted statement is carried by the author's own lemma environment in the chapter source.\par\medskip
\noindent\textit{Allowed standard prerequisites (editorial).}\begin{itemize}
\item Vector space bases: an E-basis of F and a k-basis of E; k-linear independence in E
\end{itemize}\medskip
\begin{lemma}[Multiplicativity]
\label{lemma-multiplicativity-degrees}
Suppose given a tower of fields $F/E/k$. Then
$$
[F:k] = [F:E][E:k]
$$
\end{lemma}

\begin{proof}
Let $\alpha_1, \ldots, \alpha_n \in F$ be an $E$-basis for $F$. Let
$\beta_1, \ldots, \beta_m \in E$ be a $k$-basis for $E$. Then the claim is
that the set of products
$\{\alpha_i \beta_j, 1 \leq i \leq n, 1 \leq j \leq m\}$
is a $k$-basis for $F$. Indeed, let us check first that they span $F$ over $k$.

\medskip\noindent
By assumption, the $\{\alpha_i\}$ span $F$ over $E$. So if
$f \in F$, there are $a_i \in E$ with
$$
f = \sum\nolimits_i a_i \alpha_i,
$$
and, for each $i$, we can write $a_i = \sum b_{ij} \beta_j$ for some
$b_{ij} \in k$. Putting these together, we find
$$
f = \sum\nolimits_{i,j} b_{ij} \alpha_i \beta_j,
$$
proving that the $\{\alpha_i \beta_j\}$ span $F$ over $k$.

\medskip\noindent
Suppose now that there existed a nontrivial relation
$$
\sum\nolimits_{i,j} c_{ij} \alpha_i \beta_j = 0
$$
for the $c_{ij} \in k$. In that case, we would have
$$
\sum\nolimits_i \alpha_i \left( \sum\nolimits_j c_{ij} \beta_j \right) = 0,
$$
and the inner terms lie in $E$ as the $\beta_j$ do. Now $E$-linear
independence of the $\{\alpha_i\}$ shows that the inner sums are all zero.
Then $k$-linear independence of the $\{\beta_j\}$ shows that the
$c_{ij}$ all vanish.
\end{proof}

\end{document}
